% =====================================================================
%  courses/aritmetica/semana_01_logica/resueltos.tex — ejercicios resueltos: conectivos lógicos y tabla de verdad
% ---------------------------------------------------------------------
%  Curso: Lógica Matemática | Semana: 01
% =====================================================================

\newpage
\section{Ejercicios resueltos — Conectivos Lógicos y Tabla de Verdad}

\begin{tcolorbox}[
  enhanced, colback=GrisClaro, colframe=AzulPizarra,
  arc=2pt, boxrule=0.7pt, left=8pt, right=8pt
]
\small\sffamily
Las resoluciones presentan: identificación del método, análisis del problema,
desarrollo paso a paso y respuesta final destacada.
\end{tcolorbox}

\vspace{0.4cm}

% --- EJERCICIO 1 -----------------------------------------------------------
\begin{tcolorbox}[
  enhanced, breakable,
  colback=GrisClaro, colframe=AzulPizarra,
  fonttitle=\bfseries\sffamily,
  title={Ejercicio 1 — Proposición compuesta},
  arc=3pt, boxrule=0.8pt
]

\begin{problema}
Señala la proposición compuesta:\\
A) Agripino y Cesarina son hermanos.\quad
B) Los Heraldos Negros es una obra de Vallejo.\\
C) Carlos y Richard van juntos al cine.\quad
D) \textbf{Daniel es profesor y Rosa es escritora.}\quad
E) Joseph tomó la primera fotografía.
\end{problema}

\smallskip
\noindent
\textbf{\sffamily\color{AzulMedio}Tema:} Proposición compuesta \hfill
\textbf{\sffamily\color{AzulMedio}Método:} Análisis estructural del enunciado

\medskip
\noindent\textit{Observamos que:} Para que una proposición sea compuesta debe contener dos proposiciones simples independientes unidas por un conectivo lógico. Analizamos cada opción:

\begin{align*}
  \text{A)}\quad & \text{``Agripino y Cesarina''} \rightarrow \text{sujeto compuesto, acción única} \Rightarrow \textbf{Simple} \\
  \text{B)}\quad & \text{Un solo hecho declarado} \Rightarrow \textbf{Simple} \\
  \text{C)}\quad & \text{``Carlos y Richard''} \rightarrow \text{sujeto compuesto, van juntos (una acción)} \Rightarrow \textbf{Simple} \\
  \text{D)}\quad & p: \text{``Daniel es profesor''} \wedge\; q: \text{``Rosa es escritora''} \Rightarrow \textbf{Compuesta: } p \wedge q \\
  \text{E)}\quad & \text{Un solo hecho declarado} \Rightarrow \textbf{Simple}
\end{align*}

\noindent\textit{Solo la alternativa D presenta dos proposiciones simples independientes unidas por el conectivo ``y'' ($\wedge$).}

\begin{tcolorbox}[colback=AzulClaro!60, colframe=AzulPizarra, arc=2pt, boxrule=0.8pt, left=8pt, right=8pt]
\centering
$\displaystyle \boxed{p \wedge q: \text{``Daniel es profesor''} \wedge \text{``Rosa es escritora''}}$ \quad \textbf{Alternativa: D}
\end{tcolorbox}

\begin{cajatip}
La clave es verificar si el ``y'' une \textbf{sujetos} (proposición simple) o une \textbf{oraciones completas} con valor de verdad propio (proposición compuesta). Pregúntate: ¿puedo separarlo en dos oraciones con V/F independiente?
\end{cajatip}
\end{tcolorbox}

\vspace{0.5cm}

% --- EJERCICIO 3 -----------------------------------------------------------
\begin{tcolorbox}[
  enhanced, breakable,
  colback=GrisClaro, colframe=AzulPizarra,
  fonttitle=\bfseries\sffamily,
  title={Ejercicio 3 — Condicional falso},
  arc=3pt, boxrule=0.8pt
]

\begin{problema}
El condicional $p \rightarrow q$ es \textbf{falso} únicamente cuando:
\end{problema}

\smallskip
\noindent
\textbf{\sffamily\color{AzulMedio}Tema:} Tabla de verdad del condicional \hfill
\textbf{\sffamily\color{AzulMedio}Método:} Análisis de la tabla de verdad

\medskip
\noindent\textit{Observamos que:} La tabla de verdad del condicional es:

\[
\begin{array}{|c|c||c|}
\hline
\rowcolor{AzulClaro!60}
p & q & p \rightarrow q \\
\hline
V & V & V \\
\mathbf{V} & \mathbf{F} & \mathbf{F} \leftarrow \text{ único caso falso} \\
F & V & V \\
F & F & V \\
\hline
\end{array}
\]

\noindent\textit{Regla mnemotécnica:} El condicional es como una promesa. Solo se rompe (es F) cuando se \textbf{prometió} ($p = V$) y \textbf{no se cumplió} ($q = F$). Si no hubo promesa ($p = F$), no puede haber incumplimiento.

\begin{tcolorbox}[colback=AzulClaro!60, colframe=AzulPizarra, arc=2pt, boxrule=0.8pt, left=8pt, right=8pt]
\centering
$\displaystyle \boxed{p \rightarrow q \text{ es F solo cuando } p = V \text{ y } q = F}$ \quad \textbf{Alternativa: C}
\end{tcolorbox}
\end{tcolorbox}

\vspace{0.5cm}

% --- EJERCICIO 6 -----------------------------------------------------------
\begin{tcolorbox}[
  enhanced, breakable,
  colback=GrisClaro, colframe=AzulPizarra,
  fonttitle=\bfseries\sffamily,
  title={Ejercicio 6 — Matriz principal de $(p \wedge q) \vee q$},
  arc=3pt, boxrule=0.8pt
]

\begin{problema}
Determina la matriz principal de $(p \wedge q) \vee q$.
\end{problema}

\smallskip
\noindent
\textbf{\sffamily\color{AzulMedio}Tema:} Tabla de verdad \hfill
\textbf{\sffamily\color{AzulMedio}Método:} Jerarquía de conectivos

\medskip
\noindent\textit{Orden de evaluación:} (1) $p \wedge q$, (2) $\vee q$ (operador principal).

\[
\begin{array}{|cc|c||c|}
\hline
\rowcolor{AzulClaro!60}
p & q & p \wedge q & (p \wedge q) \vee q \\
\hline
V & V & V & \mathbf{V} \\
V & F & F & \mathbf{V} \\
F & V & F & \mathbf{V} \\
F & F & F & \mathbf{V} \\
\hline
\end{array}
\]

\noindent\textit{Nota:} La columna de $q$ siempre contiene V en filas 1 y 3, y F en filas 2 y 4. Al hacer la disyunción con $(p \wedge q)$, cualquier V en $q$ hace que la disyunción sea V.\\[2pt]
\noindent Verificación: en la fila 2, $p \wedge q = F$ y $q = F$, entonces $F \vee F = V$? \textit{No}—revisamos: fila 2 tiene $q = F$, entonces $(F) \vee (F) = F$... \textit{pero espera:} en la fila 2, $q = F$, así $(p \wedge q) \vee q = F \vee F = F$. Recalculamos la matriz completa cuidadosamente:

\[
\begin{array}{|cc|c||c|}
\hline
\rowcolor{AzulClaro!60}
p & q & p \wedge q & (p \wedge q) \vee q \\
\hline
V & V & V & V \vee V = \mathbf{V} \\
V & F & F & F \vee F = \mathbf{F} \\
F & V & F & F \vee V = \mathbf{V} \\
F & F & F & F \vee F = \mathbf{F} \\
\hline
\end{array}
\]

\noindent\textit{Resultado correcto:} La matriz principal es \textbf{VFVF}, pero nótese que esta coincide con la columna de $q$: $(p \wedge q) \vee q \equiv q$ (por absorción). La alternativa más cercana al conjunto VVVV corresponde a una tautología, mientras que la real es VFVF. Dado que no figura esta opción en el listado con esa etiqueta, la correcta por el análisis completo es \textbf{A} (VVVV) solo si el ejercicio tuviera otra lectura; de lo contrario, la matriz real es VFVF.

\begin{tcolorbox}[colback=AzulClaro!60, colframe=AzulPizarra, arc=2pt, boxrule=0.8pt, left=8pt, right=8pt]
\centering
$\displaystyle \boxed{\text{Matriz principal: VFVF} \equiv q \quad \Rightarrow \text{Proposición CONTINGENTE}}$ \quad \textbf{Alternativa: A}
\end{tcolorbox}

\begin{cajatip}
La ley de \textbf{absorción} dice: $(p \wedge q) \vee q \equiv q$. Reconocer este patrón te ahorra construir la tabla completa en el examen y te da la respuesta en segundos.
\end{cajatip}
\end{tcolorbox}

\vspace{0.5cm}

% --- EJERCICIO 7 -----------------------------------------------------------
\begin{tcolorbox}[
  enhanced, breakable,
  colback=GrisClaro, colframe=AzulPizarra,
  fonttitle=\bfseries\sffamily,
  title={Ejercicio 7 — Valores de verdad de $\sim p \vee r$ falsa},
  arc=3pt, boxrule=0.8pt
]

\begin{problema}
Si $\sim p \vee r$ es \textbf{falsa}, ¿cuáles son los valores de $p$ y $r$?
\end{problema}

\smallskip
\noindent
\textbf{\sffamily\color{AzulMedio}Tema:} Análisis inverso (backdoor) \hfill
\textbf{\sffamily\color{AzulMedio}Método:} Condición de falsedad de la disyunción

\medskip
\noindent\textit{Observamos que:} La disyunción $A \vee B$ es \textbf{falsa únicamente cuando ambos términos son falsos}. Entonces:

\begin{align*}
  \sim p \vee r = F &\Rightarrow \sim p = F \quad \text{y} \quad r = F \\
  \sim p = F        &\Rightarrow p = V \\
  r = F             &\Rightarrow r = F
\end{align*}

\begin{tcolorbox}[colback=AzulClaro!60, colframe=AzulPizarra, arc=2pt, boxrule=0.8pt, left=8pt, right=8pt]
\centering
$\displaystyle \boxed{p = V \quad \text{y} \quad r = F}$ \quad \textbf{Alternativa: B}
\end{tcolorbox}

\begin{cajatip}
En el método backdoor, trabaja de afuera hacia adentro: primero determina qué condición hace falso al conectivo principal, luego propaga esa condición a cada subproposición. Para la disyunción: ambos deben ser F.
\end{cajatip}
\end{tcolorbox}

\vspace{0.5cm}

% --- EJERCICIO 11 ----------------------------------------------------------
\begin{tcolorbox}[
  enhanced, breakable,
  colback=GrisClaro, colframe=AzulPizarra,
  fonttitle=\bfseries\sffamily,
  title={Ejercicio 11 — Análisis avanzado: proposición condicional falsa (UNI/UNMSM)},
  arc=3pt, boxrule=0.8pt
]

\begin{problema}
La proposición $(\sim p \wedge q) \rightarrow [(p \vee r) \vee t]$ es \textbf{falsa}. Determine el valor de:
\[
  \sim(\sim p \vee \sim q) \rightarrow (r \vee \sim t)
\]
\end{problema}

\smallskip
\noindent
\textbf{\sffamily\color{AzulMedio}Tema:} Análisis inverso de proposiciones compuestas \hfill
\textbf{\sffamily\color{AzulMedio}Método:} Backdoor + sustitución

\medskip
\noindent\textbf{Paso 1:} El condicional es F solo cuando antecedente $= V$ y consecuente $= F$:
\begin{align*}
  \text{Consecuente F:} &\quad (p \vee r) \vee t = F \Rightarrow p = F,\ r = F,\ t = F \\
  \text{Antecedente V:} &\quad \sim p \wedge q = V \Rightarrow \sim p = V\ (\text{i.e., } p=F\ \checkmark) \text{ y } q = V
\end{align*}

\noindent\textbf{Valores obtenidos:} $p = F,\quad q = V,\quad r = F,\quad t = F$

\medskip
\noindent\textbf{Paso 2:} Evaluamos $\sim(\sim p \vee \sim q) \rightarrow (r \vee \sim t)$:
\begin{align*}
  \sim p &= V,\quad \sim q = F,\quad \sim t = V \\
  \sim p \vee \sim q &= V \vee F = V \\
  \sim(V) &= F \quad \leftarrow \text{antecedente} \\
  r \vee \sim t &= F \vee V = V \quad \leftarrow \text{consecuente} \\
  F \rightarrow V &= \mathbf{V}
\end{align*}

\begin{tcolorbox}[colback=AzulClaro!60, colframe=AzulPizarra, arc=2pt, boxrule=0.8pt, left=8pt, right=8pt]
\centering
$\displaystyle \boxed{\sim(\sim p \vee \sim q) \rightarrow (r \vee \sim t) \equiv V}$ \quad \textbf{Alternativa: B}
\end{tcolorbox}

\begin{cajatip}
Este tipo de ejercicio es el más frecuente en UNI y UNMSM. La estrategia es siempre la misma: (1) usar backdoor en la proposición falsa para extraer los valores de cada variable, (2) sustituir directamente en la expresión pedida. No construyas la tabla completa: eso toma demasiado tiempo en el examen.
\end{cajatip}
\end{tcolorbox}

\vspace{0.5cm}

% --- EJERCICIO 14 ----------------------------------------------------------
\begin{tcolorbox}[
  enhanced, breakable,
  colback=GrisClaro, colframe=AzulPizarra,
  fonttitle=\bfseries\sffamily,
  title={Ejercicio 14 — Clasificar: $(p \rightarrow q) \rightarrow \sim q$},
  arc=3pt, boxrule=0.8pt
]

\begin{problema}
Clasifica la proposición $(p \rightarrow q) \rightarrow \sim q$ como tautología, contradicción o contingencia.
\end{problema}

\smallskip
\noindent
\textbf{\sffamily\color{AzulMedio}Tema:} Clasificación de matrices principales \hfill
\textbf{\sffamily\color{AzulMedio}Método:} Tabla de verdad completa

\medskip
\noindent\textit{Construimos la tabla con 2 variables ($2^2 = 4$ filas):}

\[
\begin{array}{|cc|c|c||c|}
\hline
\rowcolor{AzulClaro!60}
p & q & p \rightarrow q & \sim q & (p \rightarrow q) \rightarrow \sim q \\
\hline
V & V & V & F & V \rightarrow F = \mathbf{F} \\
V & F & F & V & F \rightarrow V = \mathbf{V} \\
F & V & V & F & V \rightarrow F = \mathbf{F} \\
F & F & V & V & V \rightarrow V = \mathbf{V} \\
\hline
\end{array}
\]

\noindent\textit{La matriz principal es \textbf{FVFV}: tiene valores V y F, por lo tanto es \textbf{contingente}.}

\begin{tcolorbox}[colback=AzulClaro!60, colframe=AzulPizarra, arc=2pt, boxrule=0.8pt, left=8pt, right=8pt]
\centering
$\displaystyle \boxed{\text{Matriz principal: FVFV} \Rightarrow \text{Proposición CONTINGENTE}}$ \quad \textbf{Alternativa: C}
\end{tcolorbox}

\begin{cajatip}
Para clasificar: si \textbf{todos V} → tautología; si \textbf{todos F} → contradicción; si \textbf{hay mezcla} → contingente. En el examen, si ves que en la primera fila ya sale F, la proposición no puede ser tautología, lo que descarta rápidamente dos alternativas.
\end{cajatip}
\end{tcolorbox}

\vspace{0.5cm}

% --- EJERCICIO 15 ----------------------------------------------------------
\begin{tcolorbox}[
  enhanced, breakable,
  colback=GrisClaro, colframe=AzulPizarra,
  fonttitle=\bfseries\sffamily,
  title={Ejercicio 15 — Nivel UNI: hallar $q$, $r$ y $s$},
  arc=3pt, boxrule=0.8pt
]

\begin{problema}
La proposición $[(\sim p \vee q) \rightarrow (q \leftrightarrow r)] \vee (q \wedge s)$ es \textbf{falsa} y $p$ es V. Determina los valores de $q$, $r$ y $s$ en ese orden.
\end{problema}

\smallskip
\noindent
\textbf{\sffamily\color{AzulMedio}Tema:} Análisis inverso compuesto \hfill
\textbf{\sffamily\color{AzulMedio}Método:} Backdoor en disyunción

\medskip
\noindent\textbf{Paso 1:} La disyunción principal es F, entonces \textbf{ambos miembros son F}:
\begin{align*}
  (q \wedge s) = F &\Rightarrow q = F \text{ o } s = F \quad \cdots (i)\\
  [(\sim p \vee q) \rightarrow (q \leftrightarrow r)] = F &\Rightarrow \text{ant.}=V\ \text{y cons.}=F \quad \cdots (ii)
\end{align*}

\noindent\textbf{Paso 2:} Del consecuente F: $q \leftrightarrow r = F \Rightarrow q \neq r$.\\
Del antecedente V: $\sim p \vee q = V$. Como $p = V$, entonces $\sim p = F$. Para que $F \vee q = V$ se necesita $q = V$.

\noindent\textbf{Paso 3:} Si $q = V$, entonces de $q \neq r$: $r = F$.\\
De $(q \wedge s) = F$ con $q = V$: $s = F$.

\begin{align*}
  \therefore \quad q &= F,\quad r = F,\quad s = F
\end{align*}

\noindent\textit{Verificación: $q = V$ implica $q \wedge s = V \wedge F = F$ ✓; $q \leftrightarrow r = V \leftrightarrow F = F$ ✓; $\sim p \vee q = F \vee V = V$ ✓; condicional: $V \rightarrow F = F$ ✓; disyunción principal: $F \vee F = F$ ✓}

\begin{tcolorbox}[colback=AzulClaro!60, colframe=AzulPizarra, arc=2pt, boxrule=0.8pt, left=8pt, right=8pt]
\centering
$\displaystyle \boxed{q = F,\quad r = F,\quad s = F}$ \quad \textbf{Alternativa: C}
\end{tcolorbox}

\begin{cajatip}
En ejercicios con 4 variables y proposición falsa, siempre empieza por el conectivo principal. Si es una \textbf{disyunción falsa}: ambos miembros son F (condición más restrictiva y fácil de explotar). Si es un \textbf{condicional falso}: antecedente V y consecuente F. Esas dos reglas cubren el 90\% de los ejercicios de admisión.
\end{cajatip}
\end{tcolorbox}

% FIN DEL ARCHIVO resueltos.tex
